\documentclass[11pt]{article}
\usepackage[margin=1in]{geometry}
\usepackage{amsmath,amssymb,amsthm,mathtools}
\usepackage{graphicx}
\usepackage{microtype}
\usepackage{xcolor}
\usepackage[colorlinks=true,linkcolor=blue!55!black,citecolor=blue!55!black,urlcolor=blue!55!black]{hyperref}
\usepackage{enumitem}

\newtheorem{theorem}{Theorem}
\newtheorem{corollary}[theorem]{Corollary}
\newtheorem{lemma}[theorem]{Lemma}
\newcommand{\vol}{\operatorname{vol}}
\newcommand{\iq}{\operatorname{iq}}

\title{Balanced Simplex Joins Interpolate the Vertex Endpoints}
\author{Candidate partial result for arXiv:2509.13898\\Model: GPT5.6}
\date{19 July 2026}

\begin{document}
\maketitle

\begin{center}
\fbox{\parbox{0.91\textwidth}{\textbf{Status: partial result, likely valid.}
For every $1\le r\le n$ we construct an $n$-polytope with exactly $n+r$
vertices and isoperimetric quotient $\asymp n/\sqrt r$.  The matching
universal lower bound remains open.}}
\end{center}

\section{Source problem}

For a convex body $K\subset\mathbb R^n$, write
\[
  \iq(K)=\frac{\vol_{n-1}(\partial K)}{\vol_n(K)^{(n-1)/n}}.
\]
Ball, B\"or\"oczky, and Naor ask how the optimal universal lower bound on
$\iq(K)$ interpolates when an $n$-dimensional polytope has $\beta$ vertices
with $n+1<\beta<2n$ \cite[Remark~2]{BBN}.  The two endpoints are represented
by the regular simplex, with quotient of order $n$, and the cross-polytope,
with quotient of order $\sqrt n$.

\begin{figure}[ht]
  \centering
  \includegraphics[width=\textwidth]{figures/open_problem_crop.png}
  \caption{Remark 2 on page 2 of the source paper.}
\end{figure}

To state the contribution compactly, define
\[
 \mathcal I_n(\beta)=\inf\{\iq(P): P\subset\mathbb R^n
 \text{ is a convex polytope with exactly $\beta$ vertices}\}.
\]

\section{Idea of the proof}

The endpoint bodies are themselves joins of regular simplices: one
$n$-simplex at $r=1$, and the join of $n$ one-dimensional simplices at
$r=n$, which is the cross-polytope.  The natural interpolation is therefore
to join $r$ regular simplices whose dimensions are as equal as possible.

The source paper already computes the volume and surface area of an arbitrary
join (called an $\ell_1$-sum there) \cite[Lemma~12, equations (24)--(26)]{BBN}.
Scaling a $b$-simplex to inradius $1/\sqrt b$ makes the surface factor from
that lemma add exactly as $\sqrt{b_1+\cdots+b_r}=\sqrt n$.  What remains is a
single balanced product.  Stirling's formula turns that product into
$\sqrt{n/r}$, yielding $n/\sqrt r$ after the surface prefactor is restored.

\section{The partial result}

\begin{theorem}[Balanced simplex-join interpolation]\label{thm:main}
There are universal constants $c,C>0$ such that for every $n\ge1$ and every
integer $r$ with $1\le r\le n$, there exists an $n$-dimensional convex
polytope $K_{n,r}\subset\mathbb R^n$ with exactly $n+r$ vertices and
\[
  c\frac{n}{\sqrt r}\le \iq(K_{n,r})\le C\frac{n}{\sqrt r}.
\]
In particular,
\[
  \mathcal I_n(n+r)\le C\frac{n}{\sqrt r}.
\]
\end{theorem}

\begin{proof}
Choose positive integers $b_1,\ldots,b_r$ with
\[
 b_1+\cdots+b_r=n,
 \qquad b_i\in\{\lfloor n/r\rfloor,\lceil n/r\rceil\}.
\]
For each $i$, let $\Delta_i\subset\mathbb R^{b_i}$ be a regular simplex,
centered at the origin and scaled so that its inradius is
$h_i=1/\sqrt{b_i}$.  In mutually orthogonal coordinate blocks, form their
join
\[
 K_{n,r}=\operatorname{conv}(\Delta_1'\cup\cdots\cup\Delta_r')
 \subset \mathbb R^{b_1}\times\cdots\times\mathbb R^{b_r}=\mathbb R^n.
\]
This is precisely the $\ell_1$-sum in equation (24) of \cite{BBN}.  The
vertex count in Lemma 12 of that paper gives
\[
  \#\operatorname{vert}(K_{n,r})=
  \sum_{i=1}^r(b_i+1)=n+r.
\]

We next record the simplex normalization.  A regular $b$-simplex of side
length $a$ has volume
\[
 \frac{a^b\sqrt{b+1}}{b!\sqrt{2^b}}
 \quad\text{and inradius}\quad
 \frac{a}{\sqrt{2b(b+1)}}.
\]
Consequently, a regular $b$-simplex with inradius $h$ has volume
\[
 \frac{h^b b^{b/2}(b+1)^{(b+1)/2}}{b!}.
\]
At $h=1/\sqrt b$ this becomes
\begin{equation}\label{eq:simplex-volume}
 b!\,\vol_b(\Delta_i)=(b+1)^{(b+1)/2}.
\end{equation}

Let
\[
  P=\prod_{i=1}^r (b_i+1)^{(b_i+1)/2}.
\]
The volume and surface formulas (25) and (26) of \cite{BBN}, together with
\eqref{eq:simplex-volume} and $\sum_i h_i^{-2}=\sum_i b_i=n$, give
\[
 \vol_n(K_{n,r})=\frac{P}{n!},
 \qquad
 \vol_{n-1}(\partial K_{n,r})
   =\frac{\sqrt n\,P}{(n-1)!}
   =n\sqrt n\,\vol_n(K_{n,r}).
\]
It follows exactly that
\begin{equation}\label{eq:exact-iq}
 \iq(K_{n,r})
   =n\sqrt n\left(\frac{P}{n!}\right)^{1/n}.
\end{equation}

Put $m=n/r\ge1$.  Balance of the blocks implies
\[
 m\le b_i+1\le m+2\le3m,
 \qquad
 \sum_{i=1}^r(b_i+1)=n+r=n\left(1+\frac1m\right).
\]
Therefore
\[
 \sqrt m
 \le P^{1/n}
 \le (3m)^{(1+1/m)/2}
 \le 3\sqrt m,
\]
where the last constant is universal.  Stirling's formula gives
$(n!)^{1/n}\asymp n$.  Substitution into \eqref{eq:exact-iq} yields
\[
 \iq(K_{n,r})\asymp \sqrt n\,\sqrt m
 =\frac{n}{\sqrt r},
\]
as claimed.
\end{proof}

\begin{corollary}[Current bracket for the open interpolation]\label{cor:bracket}
For $1\le r\le n$,
\[
 c\sqrt n\le \mathcal I_n(n+r)\le C\frac{n}{\sqrt r}.
\]
The lower estimate is the classical Euclidean isoperimetric inequality.
\end{corollary}

\section{Verification notes}

Three checks are especially diagnostic.
\begin{enumerate}[leftmargin=2em]
\item If $r=1$, then $b_1=n$ and \eqref{eq:exact-iq} is the regular-simplex
formula in equation (11) of \cite{BBN}, asymptotic to $en$.
\item If $r=n$, then every $b_i=1$ and the join is the cross-polytope.
Equation \eqref{eq:exact-iq} becomes
$2n\sqrt n/(n!)^{1/n}\sim2e\sqrt n$, again matching equation (11).
\item The first part of Lemma 12, namely equations (24)--(26), does not
require origin symmetry.  Origin symmetry is imposed only for the later
affine-minimality formula (27), which is not used here.
\end{enumerate}

The script \texttt{code/verify\_formula.py} evaluates the logarithmic form of
\eqref{eq:exact-iq} for all $1\le r\le n\le500$.  This finite computation is
only a consistency audit; the proof above is uniform and noncomputational.

\section{Limitations and novelty check}

The theorem supplies an interpolating \emph{upper construction} for the
extremal infimum.  It does not establish the conjecturally matching statement
that every $n$-polytope with $n+r$ vertices satisfies
$\iq(K)\gtrsim n/\sqrt r$.  Thus Remark 2 remains open.

Before promotion, the run's registry, solution, attempt, and proof-gap indexes
were searched for the source id and the vertex-interpolation terminology.
Exact web/arXiv searches for ``$n+r$ vertices,'' ``isoperimetric quotient,''
``join,'' and the scale $n/\sqrt r$ found the source paper but no occurrence
of this construction.  This was a bounded search rather than a complete
MathSciNet or zbMATH audit, so the novelty assessment is provisional.  The
source uses Lemma 12 for symmetric cube blocks in the regime of at least
$2n$ vertices; it does not state the balanced-simplex application above.

\begin{thebibliography}{9}
\bibitem{BBN}
K. Ball, K. J. B\"or\"oczky, and A. Naor,
\emph{Approximate isoperimetry for convex polytopes},
arXiv:2509.13898 (2025), especially Remark 2 and Lemma 12.
\end{thebibliography}

\end{document}
